\documentclass[12pt]{article}
\usepackage[utf8]{inputenc}
\usepackage[T1]{fontenc}
\usepackage{hyperref}
\title{Sicilian Language Resources for K-12 Curriculum Development}
\author{Compiled by Claude Sonnet 5.5 (claude-sonnet-5-5) for Andr\'as Kornai}
\date{}
\begin{document}
\maketitle

\section{Introduction}
Sicilian (ISO 639-3 scn) is an Indo-European language of Sicily with several million speakers. It is generally called a dialect in Italy. This dossier lists only resources that appeared in web searches. Sicilian has no single standard spelling, so any curriculum has to choose an orthography. Teaching of Sicilian is limited and mostly framed as ``Sicilian history, literature and linguistic heritage''.

\section{Primary Educational Resources}
\begin{itemize}
\item Regional Law 9/2011 of Sicily, on teaching Sicilian: it provides for the promotion and teaching of Sicilian history, literature and linguistic heritage in schools of every level. \url{https://it.wikisource.org/wiki/Legge_regionale_Sicilia_31_maggio_2011,_n._9_-_Insegnamento_del_siciliano}
\item Regional government report on the ``Identit\`a siciliana'' project: about 300 teachers trained and over 100 school projects presented. \url{https://www.regione.sicilia.it/la-regione-informa/identita-siciliana-formati-trecento-docenti-presentati-oltre-cento-progetti-nelle-scuole}
\item Press report on one hour per week dedicated to Sicilian language and culture in schools (La Sicilia). \url{https://www.lasicilia.it/news/cronaca/1002724/scuola-dal-prossimo-anno-un-ora-la-settimana-sara-dedicata-alla-lingua-siciliana.html}
\item School circular on the programme ``Scuola e Cultura Regionale in Sicilia'', implementing Law 9/2011. \url{https://old.iccapuanapardo.edu.it/wp-content/uploads/2022/05/Circolare-322-Progetto-Scuola-e-Cultura-Regionale-in-Sicilia-per-l’attuazione-della-Legge-n.-09-2011.pdf}
\item Gaetano Cipolla, \emph{Learn Sicilian} (textbook with DVD) and \emph{Learn Sicilian II}, adult-oriented courses from Arba Sicula; mentioned in the Arba Sicula coverage. \url{https://lideamagazine.com/an-evening-of-sicilian-culture-with-gaetano-cipolla/?amp=1}
\end{itemize}

\section{Linguistic and Reference Materials}
\begin{itemize}
\item \emph{Vocabolario siciliano} (Piccitto, Tropea, Trovato; Centro di studi filologici e linguistici siciliani, 1977--2002, about 5,500 pages). \url{https://en.wikipedia.org/wiki/Vocabolario_siciliano}
\item Vocabolario Atlante Multimediale della Sicilia (VAMS) record. \url{https://iris.unipa.it/handle/10447/601873}
\item Cademia Siciliana orthographic proposal (``Proposta di normalizzazione ortografica comune della lingua siciliana''); presentation document. \url{https://cademiasiciliana.org/wp-content/uploads/Prisintazzioni-da-Cademia.pdf}
\item Cademia Siciliana overview. \url{https://en.wikipedia.org/wiki/Cademia_Siciliana}
\item Sicilian orthography overview. \url{https://en.wikipedia.org/wiki/Sicilian_orthography}
\item Sicilian language overview. \url{https://en.wikipedia.org/wiki/Sicilian_language}
\item Historical dictionaries (Biundi 1857, Mortillaro 1853, Pasqualino 1785) listed at Lexilogos. \url{https://www.lexilogos.com/sicilien_dictionnaire.htm}
\item Napizia bibliography of Sicilian resources, including J.K. Bonner, \emph{An Introduction to Sicilian Grammar}. \url{https://www.napizia.com/pages/sicilian/bibliography.shtml}
\end{itemize}

\section{Religious and Community Resources}
\begin{itemize}
\item Arba Sicula (New York), a nonprofit publishing a bilingual Sicilian/English journal and books since 1979. \url{https://lideamagazine.com/an-evening-of-sicilian-culture-with-gaetano-cipolla/?amp=1}
\item Cademia Siciliana, a nonprofit founded in 2016 to promote Sicilian. \url{https://www.balarm.it/news/magazine/come-la-criusca-ma-per-il-dialetto-la-cademia-siciliana-a-tutela-del-siciliano-19902}
\item Endangered Languages Alliance page on Sicilian. \url{https://www.elalliance.org/languages/sicilian}
\end{itemize}
No Sicilian scripture translation turned up in these searches.

\section{Digital and Audio Resources}
\begin{itemize}
\item A news item on a campaign to adopt and spread disused Sicilian words on social media. \url{https://www.balarm.it/news/adotta-una-parola-siciliana-e-diffondila-sui-social-si-va-a-caccia-del-dialetto-in-disuso-117539}
\item A news item on Telegram in Sicilian. \url{https://www.balarm.it/news/telegram-parla-siciliano-se-non-ci-credi-vai-sulle-mpustazziuna-e-cancia-lingua-116371}
\end{itemize}
The searches did not confirm the Sicilian Wikipedia or Wiktionary directly. No audio corpus was found.

\section{Curriculum Development Context}
Law 9/2011 gives a legal basis for Sicilian heritage teaching. Pilot schools and teacher training are reported, but implementation seems uneven, and no complete Sicilian-medium K-12 materials were found. A curriculum would need to settle on an orthography (Cademia Siciliana versus the Centro di studi tradition), adapt adult textbooks for children, and write new graded readers.
\end{document}